%%==========================================================
%%  CAPÍTULO 1: INTRODUCCIÓN
%%==========================================================

Este capítulo sitúa el trabajo presentando el problema que lo origina, los
objetivos que se comprometió a alcanzar, el alcance dentro del cual se movió,
la metodología que lo organizó y los resultados que produjo.

El trabajo construye una herramienta de línea de comandos que genera
proyectos backend estandarizados y permite incorporarles funcionalidades
transversales reutilizables, para un equipo de desarrollo que opera en una
entidad financiera regulada. Los cuatro capítulos centrales desarrollan cada
una de las fases de la metodología, y los dos últimos cierran el trabajo
contra sus objetivos.

\section{Definición del problema}

\subsection{Planteamiento del problema}
\label{sec:planteamiento}
% Migrado desde Anteproyecto, sección 1 (Planteamiento del Problema)
En los procesos actuales de transformación digital, las organizaciones han aumentado la demanda de desarrollo y evaluación de software con el fin de soportar sus capacidades de negocio. Esta tendencia se enmarca en el reconocimiento de la transformación digital como un impulsor fundamental de la competitividad empresarial, con cerca del 90\% de las organizaciones inmersas en estas iniciativas y un gasto global proyectado que superará los 4 billones de dólares para 2027 \citep{Jie2025}. Este crecimiento ha generado una presión constante por reducir los tiempos de entrega y aumentar la eficiencia de los equipos de desarrollo, debido a que la velocidad de salida al mercado se ha convertido en un factor diferencial para la competitividad y la satisfacción de los usuarios.

Como respuesta a estos desafíos, las metodologías ágiles se han
convertido en un enfoque clave para la gestión y el desarrollo de
proyectos \citep{Palopak2024}. Sin embargo, los equipos enfrentan
tensiones de desempeño como el equilibrio entre velocidad y calidad, lo
que evidencia la persistencia de una brecha entre la velocidad esperada
por el negocio y la capacidad real de mantenerla con calidad técnica
\citep{Armbruster2025}. El Laboratorio Digital del Banco W
\citep{UniversidadIcesi2023} es un caso de esa brecha, pues opera bajo
marcos ágiles con prácticas como Dual Track Agile \citep{MartyCagan2012}
y, a medida que logra entregar más rápido, las expectativas
institucionales aumentan y la agilidad pasa de ser una ventaja a
convertirse en el nuevo estándar exigido.

En su práctica habitual, el inicio de un proyecto implica configurar
desde cero la estructura base del sistema, las convenciones técnicas y
las integraciones con servicios externos, mientras que la incorporación
de funcionalidades transversales entre aplicaciones, como autenticación,
observabilidad o mensajería, se realiza de forma manual reutilizando
código existente. El resultado son funcionalidades duplicadas e
inconsistencias en la estructura de los proyectos que dificultan su
mantenimiento y la evolución controlada del ecosistema de aplicaciones.
Las plantillas base o \emph{boilerplates} buscan encapsular esa
estructura común, pero su reutilización real suele ser limitada porque
tienden a materializarse como artefactos estáticos que quedan
desactualizados o se acoplan a contextos específicos, y se ha señalado
que rara vez logran ser reutilizables de forma general debido a la
diversidad de necesidades y tecnologías entre proyectos
\citep{SeanPromsopeak2022}. Persiste por tanto la necesidad de una base
técnica que permita industrializar el arranque de proyectos y la
composición de capacidades reutilizables sin sacrificar la autonomía de
los equipos ni los atributos de calidad del software.

\subsection{Formulación del problema}
\label{sec:formulacion}

De lo anterior se desprende la pregunta que orienta el trabajo:

\begin{quote}
¿Cómo reducir el esfuerzo de iniciar un proyecto backend y de incorporarle
funcionalidades transversales reutilizables, en un equipo de desarrollo de
una entidad financiera regulada, sin sacrificar la autonomía del equipo sobre
el código que recibe ni los atributos de calidad del software?
\end{quote}

\subsection{Sistematización del problema}
\label{sec:sistematizacion}

La pregunta general se desagrega en cuatro preguntas guía, cada una de las
cuales se responde en un capítulo del trabajo.

\begin{enumerate}
\item ¿Qué necesita realmente un equipo de desarrollo para iniciar un
proyecto backend estandarizado, y qué parte de esa necesidad no resuelve
ninguna de las soluciones disponibles en el mercado?

\item ¿Bajo qué enfoque, qué arquitectura y qué reglas de gobierno debe
estructurarse una familia de proyectos para que un mismo activo pueda servir
a varios de ellos sin que diverjan con el tiempo?

\item ¿De qué manera puede automatizarse la generación de un proyecto
conforme a esa estructura, de modo que el resultado arranque sin
intervención manual y siga perteneciendo al equipo que lo recibe?

\item ¿Cómo puede incorporarse una funcionalidad transversal a un proyecto
ya en curso, y propagarse después un cambio de esa funcionalidad, sin
intervenir proyecto por proyecto ni sobrescribir el código propio de cada
equipo?
\end{enumerate}


\section{Objetivos del proyecto}

\subsection{Objetivo General}
% Migrado desde Anteproyecto, sección 2.1 (Objetivo General)
Desarrollar una herramienta CLI para entornos Node.js (v22–v24) que automatice la generación de proyectos backend bajo principios de arquitectura hexagonal y Domain-Driven Design, alineados con los estándares BIAN e ISO 20022, para reducir el esfuerzo de configuración al iniciar proyectos.

\subsection{Objetivos específicos}
% Migrado desde Anteproyecto, sección 2.2 (Objetivos específicos)
\begin{enumerate}
\item Definir y documentar un marco de referencia que establezca las convenciones, patrones de desarrollo y criterios de organización modular del código, así como los lineamientos de alineación con los estándares BIAN e ISO 20022 que guiarán el desarrollo de la herramienta.

\item Implementar un mecanismo de generación de proyectos mediante plantillas configurables que permita obtener una estructura base ejecutable alineada con el marco de referencia definido.

\item Construir un gestor de módulos reutilizables que permita incorporar funcionalidades transversales mediante contratos de integración definidos y compartidos entre proyectos.

\item Evaluar el comportamiento de la solución desarrollada en el contexto del Laboratorio Digital del Banco W midiendo el tiempo requerido para la generación de un proyecto frente a la línea base de proyectos existentes, la correcta incorporación de un módulo reutilizable y el inicio de implementación de lógica de negocio en el proyecto generado.

\end{enumerate}

\subsection{Divergencias declaradas respecto de los objetivos}
\label{sec:divergencias}

Dos cláusulas de los objetivos transcritos arriba no se cumplieron tal como
están formuladas, y se declaran aquí porque el lector debe saberlo antes de
leer lo que sigue. La primera es el rango de versiones de la plataforma de
ejecución, porque la herramienta exige la versión 24 y rechaza de forma
explícita las dos anteriores, en lugar de soportar el rango completo que el
objetivo enuncia. La segunda es la alineación con ISO 20022, que se revisó
como fundamento del marco de referencia y se excluyó de forma deliberada, con
la razón que la sección~\ref{sec:fundamentos-adoptados} desarrolla.

Ninguna de las dos es un descuido, y las dos tienen una decisión razonada
detrás. El Capítulo~\ref{cap:evaluacion} las recoge, junto con el resto de
las divergencias del proyecto, y da el veredicto de cada objetivo cláusula
por cláusula.

\section{Delimitaciones y alcances}
\label{sec:delimitaciones-alcances}
% Migrado desde Anteproyecto, sección 4 (Alcance del Proyecto)
% Ajuste de tiempo verbal: "se enfocará" -> "se enfocó"; "estará orientado" -> "estuvo orientado"; "se contemplará" -> "se contemplaron"; "quedarán" -> "quedaron"
El proyecto se enfocó en el diseño, implementación y validación de una herramienta CLI orientada a la creación de nuevos proyectos de software a partir de parámetros y configuraciones predefinidas. Además, estuvo orientado principalmente al desarrollo de aplicaciones backend.

Se contemplaron los siguientes componentes y capacidades:

\begin{itemize}
    \item Desarrollo de la herramienta para su uso en un entornos de desarrollo Node.js (v22-v24).
    \item Documentación para su instalación, configuración e implementación.
    \item Generación de una estructura base de proyecto que funcione como plantilla inicial.
    \item Capacidad de iniciar nuevos proyectos a partir de dichas plantillas.
    \item Mecanismo para la implementación de funcionalidades compartidas mediante reglas definidas.
    \item Desarrollo de un sistema de almacenamiento de plantillas, asegurando su disponibilidad para los equipos.
    \item Generación automática de documentación openAPI basada en los módulos y funcionalidades integradas.
\end{itemize}

Quedaron fuera del alcance:

\begin{itemize}
    \item Integración con procesos de despliegue, infraestructura o automatización DevOps.
    \item Implementación de funcionalidades específicas del negocio.
    \item Soporte para otros lenguajes de programación o tecnologías diferentes a Node.js.
\end{itemize}
Más allá de ese alcance funcional, el trabajo operó bajo tres delimitaciones
que condicionan cómo deben leerse sus resultados.

La primera es de contexto. El problema se levantó sobre un único caso de
estudio, el Laboratorio Digital del Banco W, con un equipo técnico de cuatro
desarrolladores. Las necesidades que sostienen la especificación se
contrastaron con desarrolladores de otras organizaciones, pero el contexto de
aplicación es uno solo, y las conclusiones se formulan sobre propiedades
transferibles y no sobre el caso particular.

La segunda es de confidencialidad. El trabajo se desarrolla sobre una
organización del sector financiero, de modo que no se documenta información
privada de la entidad ni del equipo. El contexto se describe por sus
propiedades, las decisiones se sostienen sobre necesidades y requerimientos
verificables, y ningún dato de identificación personal se incorpora al
documento.

La tercera es de disponibilidad del estándar. La alineación con BIAN se
sostiene sobre el documento oficial de cada capacidad de negocio, y el
proyecto dispuso de ese documento para una sola de ellas. Las demás conservan
una descripción derivada de las vistas generales del modelo, y así quedan
declaradas, porque reemplazar una paráfrasis por otra no es transcribir un
estándar.

\section{Metodología de la investigación}
% Migrado desde Anteproyecto, sección 7 (Metodología del proyecto)
% Ajuste de tiempo verbal: "se adopta" -> "se adoptó"; "orientarán/fundamentarán/conformarán" -> "orientaron/fundamentaron/conformaron"; "Diseña y construye" -> "Diseñó y construyó" (Fases 2 y 3); "Evalúa" -> "Evaluó" (Fase 4)
Con el fin de dar cumplimiento al objetivo general propuesto, se adoptó el enfoque de Design Science Research (DSR) como metodología general del proyecto \citep{Wieringa2014}. Este enfoque plantea la construcción de un artefacto que solucione de manera útil y efectiva un problema en un dominio determinado, lo cual se alinea de forma directa con la naturaleza del problema identificado y con los resultados esperados del proyecto.

\begin{figure}[H]
  \scalebox{0.54}{\usetikzlibrary{trees,positioning,arrows.meta,chains,shapes.geometric,shapes,matrix,shapes.symbols}

\definecolor{bgsolid}{RGB}{220,220,220}
\definecolor{bgsub}{RGB}{240,240,240}
\definecolor{border}{RGB}{50,50,50}


\tikzset{
  solid block/.style={
    rectangle, draw=border, thick,
    fill=bgsolid, rounded corners=6pt,
    text width=19em, align=center,
    inner sep=10pt, font=\sffamily\small
  },
  dashed block/.style={
    rectangle, draw=border, thick, dashed,
    rounded corners=6pt,
    align=center,
    inner sep=10pt, font=\sffamily\small
  },
  subnode/.style={
    rectangle, draw=border, solid, thick,
    fill=bgsub, rounded corners=5pt,
    align=center, inner sep=6pt,
    minimum width=2.4cm, font=\sffamily\small
  },
  main arrow/.style={-{Stealth[length=6pt]}, thick, draw=border},
  inner arrow/.style={-{Stealth[length=5pt]}, thick, draw=border},
  inner dashed/.style={-{Stealth[length=5pt]}, thick, draw=border, dashed},
}

\begin{tikzpicture}[node distance=0.7cm]

% FASE 1
\node[solid block] (n1) at (0,0) {
  \textbf{\normalsize Fase 1 — Fundamentación y definición}\\[5pt]
  \textbf{\small Dimensión del problema}\\
  {\footnotesize Caracterización del contexto · Stakeholders}\\
  {\footnotesize Entrevistas · Análisis cualitativo}\\
  {\footnotesize Necesidades y requerimientos}\\[5pt]
  \textbf{\small Dimensión de la solución}\\
  {\footnotesize Antecedentes y soluciones existentes}\\
  {\footnotesize Principios: DDD, Hexagonal}\\
  {\footnotesize Estándares: BIAN, ISO 20022 · Convenciones}
};

% FASE 2
\node[dashed block, below=of n1] (n2) {
  \begin{tikzpicture}[node distance=0.4cm]
    \node[minimum width=16em, minimum height=0pt, inner sep=0pt] (ruler) {};

    \node[font=\sffamily\bfseries\small, align=center,
          below=0pt of ruler] (t2)
      {Fase 2 — Mecanismo de generación de proyectos};

    \node[subnode, below=0.3cm of t2, xshift=-1.5cm] (s2a)
      {\textbf{Diseñar}\\[2pt]
       {\footnotesize Plantilla base}\\
       {\footnotesize Convenciones}};

    \node[subnode, right=0.5cm of s2a] (s2b)
      {\textbf{Prototipar}\\[2pt]
       {\footnotesize Estructura modular}\\
       {\footnotesize Entorno base}};

    \node[subnode, below=0.4cm of s2b] (s2c)
      {\textbf{Refinar}\\[2pt]
       {\footnotesize Ajustar según}\\
       {\footnotesize retroalimentación}};

    \draw[inner arrow, solid] (s2a) -- (s2b);
    \draw[inner arrow] (s2b) -- (s2c);
    \draw[inner dashed] (s2c.west) -| (s2a.south);
  \end{tikzpicture}
};

% FASE 3
\node[dashed block, below=of n2] (n3) {
  \begin{tikzpicture}[node distance=0.4cm]
    \node[minimum width=16em, minimum height=0pt, inner sep=0pt] (ruler) {};

    \node[font=\sffamily\bfseries\small, align=center,
          below=0pt of ruler] (t3)
      {Fase 3 — Gestor de módulos reutilizables};

    \node[subnode, below=0.3cm of t3, xshift=-1.5cm] (s3a)
      {\textbf{Construir}\\[2pt]
       {\footnotesize Gestor de}\\
       {\footnotesize módulos}};

    \node[subnode, right=0.5cm of s3a] (s3b)
      {\textbf{Probar}\\[2pt]
       {\footnotesize Contratos de}\\
       {\footnotesize integración}};

    \node[subnode, below=0.4cm of s3b] (s3c)
      {\textbf{Refinar}\\[2pt]
       {\footnotesize Ajustar según}\\
       {\footnotesize retroalimentación}};

    \draw[inner arrow, solid] (s3a) -- (s3b);
    \draw[inner arrow] (s3b) -- (s3c);
    \draw[inner dashed] (s3c.west) -| (s3a.south);
  \end{tikzpicture}
};

% FASE 4
\node[solid block, below=of n3] (n4) {
  \textbf{\normalsize Fase 4 — Evaluación y validación}\\[4pt]
  Caso de estudio:\\ Laboratorio Digital del Banco W \\ Tiempo de configuración, incorporación de módulo, inicio de implementación
};

\draw[main arrow] (n1) -- (n2);
\draw[main arrow] (n2) -- (n3);
\draw[main arrow] (n3) -- (n4);

\end{tikzpicture}
}
  \caption{Fases de la metodología Design Science Research aplicada al proyecto}
  \label{fig:metodologia}
\end{figure}
\pagebreak

El desarrollo se organiza en cuatro fases (Figura~\ref{fig:metodologia}), las cuales se clasifican en dos tipos según su tipo de ejecución: fases secuenciales y fases iterativas.

\begin{itemize}
  \item \textbf{Fases secuenciales:} Este tipo de ejecución responde a que sus actividades requieren consolidar resultados definidos antes de avanzar, sin necesidad de refinamiento progresivo.

  \item \textbf{Fases iterativas:} Este tipo de ejecución responde a la complejidad técnica inherente a la construcción del artefacto, donde cada iteración permite refinar progresivamente el resultado a partir de los hallazgos obtenidos en la iteración anterior.
\end{itemize}

Esta sección describe la metodología, con el propósito de cada fase y
las actividades que la componen. Su ejecución y los resultados obtenidos
se desarrollan en los capítulos siguientes, dedicando un capítulo a cada
fase: la Fase 1 en el Capítulo~\ref{cap:fundamentacion}, la Fase 2 en el
Capítulo~\ref{cap:generacion}, la Fase 3 en el Capítulo~\ref{cap:gestor}
y la Fase 4 en el Capítulo~\ref{cap:evaluacion}.

\subsubsection*{Fase 1: Fundamentación y definición del marco de referencia}

Estableció las bases conceptuales que orientaron el diseño y construcción del artefacto, mediante el análisis del problema, la revisión de soluciones existentes y la definición de los estándares, convenciones y principios que estructuran la solución.

La fase se abordó en dos dimensiones complementarias. La \textbf{dimensión del problema} (actividades 1 a 4) comprendió un proceso formal de obtención de requisitos sobre un caso de estudio real, con el fin de sustentar en evidencia qué debe resolver la herramienta y bajo qué condiciones sería adoptada, en lugar de asumirlo de antemano. La \textbf{dimensión de la solución} (actividades 5 a 9) comprendió la revisión de las soluciones existentes y la definición del marco de referencia técnico que estructura el artefacto, una vez comprendido el problema.

\paragraph*{Actividades relacionadas}

\begin{enumerate}
    \item Caracterizar el contexto de aplicación en cuanto a su forma de trabajo, su contexto tecnológico y sus prácticas de reutilización de código.

    \item Identificar y justificar los stakeholders del artefacto, y definir el protocolo de entrevistas con guiones diferenciados por rol.

    \item Analizar cualitativamente las entrevistas para identificar los temas recurrentes entre roles y agruparlos en los grupos temáticos que la solución debe atender.

    \item Especificar los requerimientos funcionales y no funcionales de la herramienta, priorizados y trazables a los temas identificados.

    \item Revisar la literatura y antecedentes sobre estandarización y generación de proyectos de software.

    \item Analizar y comparar soluciones existentes en el ámbito de la generación y estandarización de proyectos de software, estableciendo un marco de evaluación e identificando las brechas que ninguna de ellas resuelve.

    \item Identificar y seleccionar los principios y buenas prácticas que fundamentaron la solución, incluyendo arquitectura hexagonal y Domain-Driven Design.

    \item Revisar y adoptar los estándares del sector financiero aplicables, como BIAN e ISO 20022, como referencia para la definición de dominios y capacidades bancarias reutilizables.

    \item Definir las convenciones técnicas y patrones de desarrollo que conformaron la herramienta.
\end{enumerate}

\subsubsection*{Fase 2: Diseño y construcción del mecanismo de generación de proyectos}

Diseñó y construyó de forma iterativa el mecanismo que permite generar proyectos con una estructura modular estandarizada desde el inicio, alineada con el marco de referencia definido en la fase anterior, reduciendo el tiempo de arranque y garantizando coherencia entre proyectos.

\paragraph*{Actividades relacionadas}

\begin{enumerate}
    \item Definir los requisitos y parámetros del mecanismo de generación de proyectos.

    \item Diseñar la estructura modular base alineada con los principios arquitectónicos y estándares bancarios definidos en el marco de referencia.

    \item Construir el mecanismo de generación de proyectos alineado con el marco de referencia definido.

    \item Configurar el entorno de desarrollo de forma reproducible, incluyendo convenciones de código y herramientas base.

    \item Implementar la generación automática de documentación openAPI basada en la estructura del proyecto generado.
\end{enumerate}

\subsubsection*{Fase 3: Diseño y construcción del gestor de modulos reutilizables}

Diseñó y construyó de forma iterativa el sistema que permite incorporar funcionalidades transversales en los proyectos de forma desacoplada, mediante contratos de integración definidos y compartidos entre proyectos, garantizando su reutilización sin generar dependencias entre su implementación concreta y la lógica de negocio de cada proyecto.

\paragraph*{Actividades relacionadas}

\begin{enumerate}
    \item Identificar y categorizar las funcionalidades transversales reutilizables entre proyectos, tomando como referencia los dominios funcionales definidos por BIAN.

    \item Definir los contratos de integración y las reglas de composición que garantizan el desacoplamiento entre módulos y proyectos.

    \item Diseñar e implementar el sistema de composición modular alineado con los contratos de integración definidos y el marco de referencia establecido.

    \item Validar que el sistema de composición modular permite incorporar funcionalidades transversales de forma desacoplada en los proyectos generados.

    \item Verificar que las funcionalidades transversales implementadas cumplen con los lineamientos de alineación con los estándares BIAN e ISO 20022 definidos en el marco de referencia.
\end{enumerate}

\subsubsection*{Fase 4: Evaluación y validación del artefacto}

Evaluó la solución desarrollada mediante un caso de estudio en el Laboratorio Digital del Banco W, midiendo el tiempo de configuración inicial frente a proyectos existentes como línea base, la correcta incorporación de un módulo reutilizable y el inicio de implementación de lógica de negocio en el proyecto generado.

\paragraph*{Actividades relacionadas}

\begin{enumerate}
    \item Definir los criterios de evaluación del caso de estudio, incluyendo la medición del tiempo de configuración inicial frente a la línea base de proyectos existentes, la incorporación de un módulo reutilizable y el inicio de implementación de lógica de negocio en el proyecto generado.

    \item Aplicar la solución en un caso de estudio en el Laboratorio Digital del Banco W, generando un proyecto base e incorporando un módulo reutilizable en un entorno de desarrollo real.

    \item Recolectar y analizar los datos obtenidos sobre el tiempo de configuración inicial, la incorporación del módulo reutilizable y el inicio de implementación de lógica de negocio, comparando los resultados frente a la línea base establecida.

    \item Elaborar el informe de evaluación con los hallazgos y recomendaciones sobre la pertinencia del artefacto.
\end{enumerate}

\section{Resultados obtenidos}

El trabajo produjo cuatro resultados, uno por fase, y cada uno se desarrolla
en el capítulo correspondiente.

\begin{enumerate}
\item \textbf{Un marco de referencia técnico} que encuadra la solución como
línea de producto de software y fija la arquitectura de la familia de
proyectos que la herramienta genera, con sus fundamentos adoptados, la
organización del código, los contratos de integración y el modelo de
gobierno de los activos. Es el insumo del que depende todo lo demás, y se
desarrolla en el Capítulo~\ref{cap:fundamentacion}.

\item \textbf{Un mecanismo de generación de proyectos}, materializado en una
herramienta de línea de comandos instalable que produce un proyecto backend
ejecutable sin intervención manual, con las convenciones del marco aplicadas,
el entorno de trabajo configurado y su documentación de API generada bajo el
estándar OpenAPI. Junto a él, un almacén de plantillas en el que una
plantilla es un activo que se describe a sí mismo, con dos plantillas
publicadas que derivan formas de producto distintas. Se desarrolla en el
Capítulo~\ref{cap:generacion}.

\item \textbf{Un gestor de módulos reutilizables} que incorpora una
funcionalidad transversal a un proyecto ya en curso, informa de su estado, la
actualiza preservando el código propio del equipo y permite retirarla, todo a
través de contratos de integración declarados en un artefacto neutral. El
catálogo cuenta con dos módulos construidos, de naturaleza distinta entre sí.
Se desarrolla en el Capítulo~\ref{cap:gestor}.

\item \textbf{Un informe de evaluación} del artefacto construido, medido con
las personas del contexto de aplicación y contra una línea base establecida
antes de que la herramienta existiera. Se desarrolla en el
Capítulo~\ref{cap:evaluacion}.
\end{enumerate}

Los activos de la línea quedaron publicados en el registro de paquetes de la
organización, que es la condición sin la cual ninguno de los tres primeros
resultados sería utilizable por un equipo distinto del que los construyó.

%%==========================================================
%%  RIESGOS ÉTICOS
%%  Se incluye como sección final del Capítulo 1.
%%==========================================================

\section{Riesgos éticos}
\label{sec:riesgos-eticos}

Un sistema que estandariza y automatiza el inicio de aplicaciones concentra
en un único punto decisiones técnicas que después se replican en varios
proyectos. En un contexto financiero esa concentración importa más, por el
manejo de información sensible y de lógica de negocio regulada. El proyecto
identificó cuatro riesgos éticos al plantearse, cada uno con la mitigación a
la que se comprometía, y lo que corresponde decir aquí es cuál de las cuatro
se cumplió. El veredicto final de cada uno, con la evidencia de la
evaluación, lo da el Capítulo~\ref{cap:evaluacion}.

\paragraph*{Exposición de información sensible en plantillas y módulos}

El compromiso era separar de forma estricta la estructura base de los valores
de configuración sensibles, promoviendo variables de entorno y mecanismos de
gestión de secretos desde el diseño. \textbf{Se cumplió, y quedó verificado
por ejecución.} Es lo que la especificación de requerimientos exige como
atributo de configurabilidad y lo que el artefacto hace, porque ninguna
dirección ni credencial vive en el código del producto, toda
configuración entra por variable de entorno validada contra un esquema, el
circuito de verificación barre el código del producto derivado buscando
valores fijos y direcciones, y el registro de eventos del producto redacta
por nombre de clave todo secreto o identificador de persona. La verificación
automática es más de lo que el compromiso pedía.

\paragraph*{Propagación de lógica financiera mal interpretada}

El compromiso era someter a revisión técnica especializada los módulos con
lógica financiera antes de incorporarlos al catálogo. \textbf{El riesgo no se
materializó}, porque ninguno de los dos módulos construidos lleva lógica
financiera, ya que uno guarda y lee registros y el otro exporta señales de
observabilidad. Las dos capacidades son técnicas y transversales por decisión
del marco, que reparte lo técnico al módulo y el dominio al producto. Lo que
queda escrito es la condición de reactivación, porque el día que el catálogo
reciba un módulo que interprete un estándar financiero, y el de mensajería
excluido en la sección~\ref{sec:fundamentos-adoptados} es el caso obvio, la
revisión especializada es exigible antes de su entrada.

\paragraph*{Exclusión de equipos con menor capacidad técnica}

El compromiso tenía dos partes y \textbf{se cumplió a medias}. La primera era
elaborar documentación accesible orientada a distintos niveles de
experiencia, y \textbf{no se construyó}. La propuesta existió, y se descartó
con un argumento de oportunidad, que esa documentación crece con los módulos
incorporados y hacerla antes documenta un producto que todavía va a cambiar
de forma. El argumento es razonable como decisión de alcance de una fase y no
es una respuesta al riesgo ético, porque nadie advirtió que lo fuera, ya que
el rechazo comparó la propuesta contra el calendario y no contra el
compromiso. Se declara así, y la documentación por niveles queda como
recomendación de adopción.

La segunda parte sí está. El producto derivado nace con documentación que
explica qué es, qué hace cada operación, cuál lleva efecto de negocio y cuál
es un andamio; el recorrido de creación se lee en lenguaje de negocio con el
nombre del estándar al lado; la herramienta declara su alcance antes de
escribir y pide confirmación; y los mensajes de error dicen qué encontró, qué
esperaba y qué resuelve la situación. Todo eso baja la barrera de adopción,
que es lo que el riesgo persigue. El acompañamiento con procesos de formación
y transferencia de conocimiento es organizacional y queda fuera de lo que un
artefacto puede resolver.

\paragraph*{Concentración de poder técnico en quienes controlan el sistema}

El compromiso era definir un modelo de gobernanza distribuido con procesos
participativos, conocimiento documentado de forma exhaustiva y una evolución
que responda a necesidades colectivas. \textbf{El mecanismo existe y el
ejercicio no}, y la razón es una condición del entorno declarada desde la
segunda fase.

El mecanismo está completo, porque el catálogo declara un responsable por
activo, la aprobación del gobierno de la línea es obligatoria sobre la ruta
del manifiesto, y el procedimiento de incorporación tiene sus pasos y sus
condiciones de validez verificables por tipo de activo. La documentación
exhaustiva es la parte más fuerte del compromiso, con casi doscientas
decisiones registradas con su contexto, su alternativa y su traza, más de
treinta documentos de expediente y un marco con su propio mecanismo de
evolución, que es lo que permite que otra persona retome el sistema. Y la
evolución respondió a necesidades que no son del autor, porque las doce
necesidades y las tres brechas salen de las entrevistas con las personas del
contexto, y varias rondas de mejora las abrieron una revisión externa y
reportes de uso.

Lo que no hay es más de un participante. El proyecto se construyó con una
sola persona en los tres roles, de modo que proponente y revisor coinciden, y
una gobernanza distribuida instanciada con un participante es un mecanismo
probado y un ejercicio ausente. Trasladarlo a operación consiste en cambiar
quién figura como responsable de cada activo, lo que el modelo de gobierno ya
contempla, y queda como recomendación de adopción.


\section{Resumen del capítulo}

Este capítulo estableció el problema, los objetivos y el marco dentro del
cual se movió el trabajo. El problema es que iniciar un proyecto backend y
dotarlo de funcionalidades transversales es, en el contexto estudiado, un
proceso manual que se concentra en pocas personas y que produce proyectos
que divergen entre sí. La respuesta comprometida es una herramienta que
genere proyectos estandarizados y permita componerles capacidades
reutilizables, construida bajo Design Science Research en cuatro fases.

Dos cosas quedaron declaradas aquí, y las dos deben retenerse al leer lo que
sigue. La primera son las dos divergencias respecto de los objetivos
comprometidos, el piso de la plataforma de ejecución y la exclusión de uno de
los dos estándares financieros, que se nombran antes de desarrollarlas en
lugar de dejarlas para el cierre. La segunda son las delimitaciones bajo las
que el trabajo operó, en particular que el problema se levantó sobre un único
caso de estudio y que la alineación con el estándar de negocio se sostiene
sobre el documento oficial de una sola capacidad.

Los cuatro capítulos siguientes desarrollan una fase cada uno, en el orden en
que se ejecutaron, y el séptimo cierra el trabajo contra los objetivos que
este capítulo transcribe. El Capítulo~\ref{cap:marco} fija antes el
vocabulario y revisa lo que existe.
